% 編輯說明（書末附錄之首）。手寫，不由腳本生成。
\apxsection{apx-intro}{編輯說明}

一、本冊重排清李漁《閒情偶寄》十六卷本之第一冊，收余懷序、凡例七則、總目、卷一目次及卷一詞曲部結構第一、詞采第二。
底本為北京大學圖書館藏康熙間翼聖堂刻本\allowbreak（Internet Archive\allowbreak\ 02097225.cn，黑白掃描）；
參校本為哈佛燕京圖書館藏同版印本\allowbreak（HOLLIS\allowbreak\ 990080074400203941，彩色掃描）。兩本出於同一副版片，裝訂次第不同
（北大本總目在凡例之前），本版從哈佛本作序、凡例、總目；書衣、扉葉及其前後空白葉之次第亦從哈佛本。
哈佛本另有扉葉「笠翁祕書第一種」及序末、凡例末印章。

扉葉北大本缺，據哈佛本重排；其中隸書書名四字，取自哈佛燕京圖書館藏本影像（哈佛圖書館對公有領域作品之數字影像不主張著作權）。

二、重排依原刻行款：一頁當原刻半葉，每行字數、低格、抬頭、雙行小注悉依原刻，版心標卷次葉碼，以便與原書對檢。

三、文字以底本為準，凡底本漫漶者據參校本定字。兩本正文除一二筆畫之別（見\apxref{apx-jiaoji}{校記}）外未見刻字之異；
北大本卷一第六、七、八、十四、十七、十八、二十六葉之眉批，哈佛本版框上空白，當為印次不同。

四、原刻異體字，依編者所定，統一改為通行字（見\apxref{apx-variants}{「異體字統一表」}）；
「于」「后」「迺」「叶」等與通行字非一詞者，及「癡」「纔」「裏」「秘」「祕」，悉照原刻。
本書「曰」「日」刻形幾無分別，一依文意定字，不出校。

五、圈點：序、總目、凡例依原刻錄出並經校對。卷一原刻圈點繁密（句讀圈、密圈、實心點、頓點），
未及逐字覆核，為便閱讀，暫不照排，改依維基文庫標點本（CC BY-SA 4.0）斷句：句末（。？！；）作小圈，句中（，、：）作頓點。
原刻圈點之錄文仍存於底稿，俟日後校定。
眉批照北大本錄出，凡不能確定之字一律作「□」；北大本有而哈佛本無之眉批十條，尤多闕字。
以上俟日後以他本校補，亦歡迎讀者指正。

六、錄文、校對與排版由 Claude（Anthropic）輔助完成，尚待通曉古籍者覆核。
